
% We care about posterior expectations, and these expectations have a
% frequentist variance.
% Even Bayesians should care about frequentist variances as a robustness measure.
Suppose we want to use Bayesian methods to infer some parameter or quantity of
interest from exchangeable data. The Bayesian posterior is often summarized by a
point estimate and measure of subjective uncertainty, which we will respectively
take to be the posterior mean and posterior variance.  The computed posterior
mean is a function of the observed data, and so has a well-defined frequentist
variance under re-sampling of the exchangeable data.  The Bayesian posterior
variance of a quantity and the frequentist variance of its mean are conceptually
distinct: the posterior variance can be thought of as describing the
uncertainty of subjective beliefs under the assumption of correct model
specification, and the frequentist variance captures the sensitivity of the
Bayesian point estimate to re-sampling of the observed data.  In general,
both may be of interest.
Even a strict Bayesian can meaningfully interpret the
frequentist variance of the posterior mean as a \emph{sensitivity analysis},
quantifying how much the Bayesian point estimate would vary had we received
different but realistic data.  Such a sensitivity analysis is particularly
germane when the data is known to have come from a random sample from a real
population. 


% You might think the two notions are the same, but they are not when you
% have misspecification
In general, the frequentist and Bayesian variances are not only conceptually
distinct but can be numerically very different.  It is true that, with
independent and identically distributed (IID) data, some reasonable regularity
conditions, a correctly specified model, and a parameter vector of fixed
dimension, the frequentist and Bayesian variances coincide asymptotically as a
consequence of the Bernstein--von Mises theorem, a.k.a.\ the Bayesian central limit
theorem.  However, under \emph{model misspecification}, the equality between the
frequentist and Bayesian variances no longer holds in general, even
asymptotically \citep{kleijn:2012:bvm} --- and most models are misspecified.


% Standard tools are inadequate
When posterior means are estimated using Markov Chain Monte Carlo (MCMC), the
two standard tools for estimating the frequentist variability of Bayesian
posterior expectations --- the nonparametric bootstrap  (henceforth simply
``bootstrap'') and the Laplace approximation --- can be each be problematic,
though for different reasons.
%
% The bootstrap is inadequate
To compute a bootstrap estimate of
the frequentist variance of a posterior mean, one repeatedly draws ``bootstrap
datasets,'' with replacement from the original set of datapoints. For each
bootstrap dataset, one then computes the posterior mean for each set of
re-sampled data \citep{huggins:2022:bayesbag}.  Though straightforward to
implement and naively parallelizable, the bootstrap remains
computationally expensive, since a separate posterior estimate (e.g., a
run of MCMC) must be computed for each bootstrap dataset.
%
% The Laplace approximation is inadequate
Alternatively, it is straightforward to estimate frequentist variance of a
Laplace approximation to the posterior.   Specifically, one can do so by
computing the ``sandwich covariance'' estimate of the maximum \textit{a
posteriori} (MAP) estimate at which the Laplace approximation is centered
\citep{stefanski:2002:mestimation}.  This approach essentially uses the standard
consistent estimator of the limiting covariance derived by
\citet{kleijn:2012:bvm}.  However, this approach requires using a particular Laplace
approximation to the posterior rather than MCMC and, as we show below, the
Laplace approximation can be practically and conceptually problematic.  For
example, unlike true posterior draws, the Laplace approximation is not invariant
to reparameterization, and, for some parameterizations the Laplace approximation
can fail catastrophically, especially in hierarchical models.


% We solve these problems with the IJ covariance.
In the present paper, we provide an estimator of the frequentist variability of
posterior expectations that enjoys the model-agnostic properties of the
bootstrap, but can be computed from a single run of MCMC. Our
method is based on the \emph{infinitesimal jackknife} (IJ) variance estimator, a
classical frequentist variance estimator derived from the \emph{empirical
influence function} of a statistical functional.  We show that the Bayesian
empirical influence function is easily computable as a particular set of
posterior covariances, the variance of which is a consistent estimator of the
frequentist variance of a posterior expectation when a Bayesian central limit
theorem (BCLT) holds for all model parameters.

% We potentially solve other problems.
Finally, though our present concern is frequentist variance estimates, the
empirical influence function is a powerful statistical tool, with applications
in robustness, cross validation, and bias estimation. Consequently, our results
for the influence function of Bayesian posteriors and a BCLT for expectations of
data-dependent functions are of broader interest.  We leave off exploring these
topics for future work.
